% Options for packages loaded elsewhere
\PassOptionsToPackage{unicode}{hyperref}
\PassOptionsToPackage{hyphens}{url}
\documentclass[
]{article}
\usepackage{xcolor}
\usepackage{amsmath,amssymb}
\setcounter{secnumdepth}{-\maxdimen} % remove section numbering
\usepackage{iftex}
\ifPDFTeX
  \usepackage[T1]{fontenc}
  \usepackage[utf8]{inputenc}
  \usepackage{textcomp} % provide euro and other symbols
\else % if luatex or xetex
  \usepackage{unicode-math} % this also loads fontspec
  \defaultfontfeatures{Scale=MatchLowercase}
  \defaultfontfeatures[\rmfamily]{Ligatures=TeX,Scale=1}
\fi
\usepackage{lmodern}
\ifPDFTeX\else
  % xetex/luatex font selection
\fi
% Use upquote if available, for straight quotes in verbatim environments
\IfFileExists{upquote.sty}{\usepackage{upquote}}{}
\IfFileExists{microtype.sty}{% use microtype if available
  \usepackage[]{microtype}
  \UseMicrotypeSet[protrusion]{basicmath} % disable protrusion for tt fonts
}{}
\makeatletter
\@ifundefined{KOMAClassName}{% if non-KOMA class
  \IfFileExists{parskip.sty}{%
    \usepackage{parskip}
  }{% else
    \setlength{\parindent}{0pt}
    \setlength{\parskip}{6pt plus 2pt minus 1pt}}
}{% if KOMA class
  \KOMAoptions{parskip=half}}
\makeatother
\usepackage{longtable,booktabs,array}
\newcounter{none} % for unnumbered tables
\usepackage{calc} % for calculating minipage widths
% Correct order of tables after \paragraph or \subparagraph
\usepackage{etoolbox}
\makeatletter
\patchcmd\longtable{\par}{\if@noskipsec\mbox{}\fi\par}{}{}
\makeatother
% Allow footnotes in longtable head/foot
\IfFileExists{footnotehyper.sty}{\usepackage{footnotehyper}}{\usepackage{footnote}}
\makesavenoteenv{longtable}
\setlength{\emergencystretch}{3em} % prevent overfull lines
\providecommand{\tightlist}{%
  \setlength{\itemsep}{0pt}\setlength{\parskip}{0pt}}
\usepackage{bookmark}
\IfFileExists{xurl.sty}{\usepackage{xurl}}{} % add URL line breaks if available
\urlstyle{same}
% fallback for those not using the hyperref driver hyperxmp:
\makeatletter
\@ifundefined{xmpquote}{\newcommand{\xmpquote}[1]{#1}}{}
\makeatother
\hypersetup{
  pdftitle={Civilodynamics --- Foundational Charter},
  pdfauthor={Satchel Hamilton · Civilodynamics Initiative},
  hidelinks,
  pdfcreator={LaTeX via pandoc}}

\title{Civilodynamics --- Foundational Charter}
\author{Satchel Hamilton · Civilodynamics Initiative}
\date{Canonical Corpus, Document 1 of 8 · Version 2.0 · June 2026}

\begin{document}
\maketitle

\begin{quote}
\textbf{Status and role.} This is the constitution of the discipline: it
states why Civilodynamics exists, what it commits to, how it is
organized, and what would prove it wrong. It supersedes the v1.0
Foundational Charter. It is deliberately \emph{not} a technical document
--- all mathematics is specified in \textbf{Document 2 (Mathematical \&
Methodological Foundations)}, to which this charter defers throughout;
the empirical application to the United States, which appeared in v1.0,
now lives in \textbf{Document 5}; and the detailed, living
research-and-dissemination program lives in \textbf{Document 8 (Program
\& Milestone Map)}. This charter sets the principles those documents
operate under.
\end{quote}

\begin{center}\rule{0.5\linewidth}{0.5pt}\end{center}

\begin{quote}
\emph{Civilodynamics is the quantitative science of how civilizations
maintain coherence, accumulate stress, approach phase transitions, and
reorganize --- with particular attention to the role of information
architecture in accelerating or damping these dynamics.}
\end{quote}

\begin{center}\rule{0.5\linewidth}{0.5pt}\end{center}

\subsection{1. The necessity of a new
discipline}\label{the-necessity-of-a-new-discipline}

The twenty-first century has made the limits of the existing social
sciences difficult to ignore. Within a single generation the global
financial system approached collapse with little warning from its
professional stewards; a series of consolidated democracies underwent
political transformations their own institutional analysts rated as
improbable; a pandemic exposed the brittleness of supply chains,
public-health systems, and epistemic institutions at once; and
algorithmic information systems reshaped the political psychology of
whole populations faster than any regulatory or scholarly framework
could track.

In retrospect, none of these were structurally unforeseeable. The
conditions that produced them were, in each case, measurable in advance.
What was missing was not data but a framework capable of integrating
that data into a coherent picture of a civilization's structural state
and trajectory. Economics studied the financial system near equilibrium.
Political science studied elections and institutions at the level of
events. Sociology described social structure without modeling its
dynamics. History narrated the past without formalizing its
regularities. Each illuminated a part while the whole moved toward
crisis unobserved.

Civilodynamics exists to address that gap. It does not aspire to predict
specific events --- the ambition of point-prediction is precisely the
failure mode of earlier grand theories. It aspires to something more
modest and more achievable: to characterize the structural state of a
civilization quantitatively, to identify when that state has entered a
regime in which qualitative disruption becomes probable, and to track
the evolution of that probability with calibrated rigor. In the language
of meteorology, Civilodynamics does not forecast the specific storm; it
measures when and where atmospheric conditions have become
storm-generative.

\subsection{2. What Civilodynamics is}\label{what-civilodynamics-is}

Four commitments distinguish the discipline and define its identity.
Each is binding; departing from any of them would make the work
something other than Civilodynamics.

\textbf{The civilization is the unit of analysis.} The object of study
is the dynamic structural state of the society as a whole --- not its
events, actors, or institutions in isolation, but the coupling between
its subsystems and the emergent behavior that coupling produces.

\textbf{The information and epistemic layer is a first-class state
variable.} The structural transformation of the information environment
is treated not as background context but as a primary dynamic variable
that alters the transfer function between structural stress and
political outcome. This is the discipline's most distinctive commitment
(§3, §5).

\textbf{Falsifiable probabilistic prediction is the standard of
validity.} Following the superforecasting tradition, claims are
expressed as probabilities over defined outcomes with stated resolution
criteria, timestamped in advance, and scored. A discipline that cannot
specify the conditions under which its claims would fail is not a
science (§6).

\textbf{Phase transitions are the central predictive object.} The
discipline targets moments of qualitative civilizational reorganization
--- the approach to bifurcation --- rather than steady-state
description. It predicts \emph{regime type}, not specific events.

\subsubsection{Scope: what the discipline does and does not
claim}\label{scope-what-the-discipline-does-and-does-not-claim}

These limits are stated up front, because honesty about them is itself a
methodological commitment (the full statement is in Document 2 §10):

\begin{itemize}
\tightlist
\item
  It \textbf{does not predict specific events} --- which war, election,
  or crisis, or when. It locates structural boundaries and estimates
  stability margins.
\item
  It \textbf{does not eliminate contingency.} Determinism in the
  underlying dynamics coexists with genuine stochasticity in the shocks
  that perturb them. (This is the formal content of the \emph{Mule
  Problem}, after Asimov's singular disruptor: the framework predicts
  that a Mule-type emergence becomes \emph{likely} in sufficiently
  stressed states, never \emph{which individual} fills the role.)
\item
  It is, in its present form, \textbf{anthropocentric} --- built on
  human cognition, metabolism, and ecology. Generalization beyond the
  human case is a research program kept deliberately separate from the
  empirical core.
\item
  Its validation is, at this stage, \textbf{promissory.} No index has
  yet been computed end-to-end against data; no retrodiction has been
  scored. The framework is a coherent, falsifiable skeleton, and the
  calibration and validation that would earn it the name of science are
  the work ahead, not a footnote to it.
\end{itemize}

\subsection{3. Intellectual lineage}\label{intellectual-lineage}

Civilodynamics does not emerge from nothing. It stands on a century and
a half of serious attempts to bring scientific rigor to the study of
human societies. Each captured something real and fell short in a
characteristic way; understanding those shortfalls is the argument for
what this discipline must do differently.

\textbf{Auguste Comte} named sociology with the explicit ambition of a
``social physics.'' The ambition was right; the execution lacked any
mathematical formalism, measurement program, or falsifiability standard.
\textbf{Marx} built a structural theory in which material conditions
constrain ideological possibility --- a lasting contribution --- but
rendered it teleological and deterministic, unable to accommodate
contingency or the reflexive modification of outcomes by actors aware of
the theory. \textbf{Talcott Parsons} achieved an abstraction so general
it could describe everything and predict nothing. \textbf{E. O.
Wilson's} sociobiology grounded individual and small-group behavior in
evolution but could not bridge to civilization-scale emergent dynamics.
The \textbf{Santa Fe Institute} complexity tradition supplied the right
conceptual repertoire --- self-organized criticality, emergence, power
laws, phase transitions --- but largely as analogy rather than literal
structural description, without a unified measurement program at
civilizational scale. \textbf{Sociophysics} (Galam, Stauffer, and
successors) applied statistical-mechanics models with exactly the right
tools, but mostly at the micro-to-meso scale of opinion dynamics.
\textbf{Mainstream economics} achieved genuine predictive success in
narrow market contexts on equilibrium assumptions that exclude, by
construction, the non-equilibrium dynamics most relevant to transitions.
\textbf{Political science} produced sophisticated event-level and
comparative work while systematically failing to anticipate major
structural disruptions, because the conditions that make a transition
likely are largely independent of the specific event that triggers it.
\textbf{Tetlock's} superforecasting established that calibrated
probabilistic reasoning can be scored and can outperform expert
intuition --- the most epistemically serious tradition in the social
sciences, and one Civilodynamics adopts wholesale --- but it is a
methodology without a theory, operating on short horizons and narrow
questions.

\textbf{Cliodynamics and the work of Peter Turchin occupy a different
place in this lineage, and the distinction must be stated plainly.}
Cliodynamics is not one predecessor among many; it is the direct
intellectual parent of Civilodynamics. The structural-demographic model,
secular-cycle theory, the practice of timestamped falsifiable
prediction, and the Seshat Global History Databank are genuine
scientific achievements, and the quantitative-historical approach this
discipline depends on would lack its empirical grounding without them.
\textbf{Civilodynamics extends cliodynamics; it does not supersede or
replace it, and it positions itself as a collaborator with that program
rather than a competitor to it.} Its claim to be a distinct discipline
rests on three --- and only three --- narrow extensions:

\begin{enumerate}
\def\labelenumi{\arabic{enumi}.}
\tightlist
\item
  \textbf{The epistemic/information layer as a first-class state
  variable.} Cliodynamics models the information environment thinly;
  Civilodynamics treats it as a primary axis of the state space and as a
  structural accelerant of the translation from stress to outcome.
\item
  \textbf{The integration of recurrent social formations with
  theoretical morphospace.} Civilodynamics explicitly identifies the
  empirical attractor structure found in the comparative record
  (Peregrine; the Seshat clustering of Shults and colleagues) with the
  morphospace tradition in evolutionary biology, yielding the
  \textbf{poliform} framework (Document 3).
\item
  \textbf{Real-time measurement.} Cliodynamics is calibrated primarily
  on historical data; Civilodynamics aims to fuse the historical
  baseline with contemporary, real-time measurement streams.
\end{enumerate}

Everything else of substance the two share. Overstating the difference
would be both inaccurate and strategically self-defeating; the
discipline's posture toward Turchin and his collaborators is one of
acknowledged debt and active alliance.

\subsection{4. The theoretical core (in
brief)}\label{the-theoretical-core-in-brief}

The formal apparatus is the subject of Document 2; what follows is a
conceptual précis sufficient to orient the reader. Nothing here is
derived --- Document 2 derives it, and where this charter and that
document might appear to differ, Document 2 governs.

\textbf{Two state descriptions.} A society is described at two coupled
levels: a slow-moving \textbf{structural configuration vector} that
locates it in the poliform morphospace along six constraint axes
(energy, coordination, information, extraction, trust, ecology), and a
faster-moving \textbf{stress readout} (economic, elite, legitimacy,
cohesion, epistemic, demographic) that measures its distance from its
own stable regime. The stress readout is a derived projection of the
configuration and its rate of change relative to a defended baseline.

\textbf{Poliforms as attractors.} Human societies do not distribute
evenly across the space of conceivable configurations; they cluster into
a small library of recurrent forms. Civilodynamics interprets these as
\textbf{basins of attraction} --- the conceptual keystone that
distinguishes the discipline from descriptive historiography. The
relationship to biology is one of \textbf{structural homology}, a shared
mathematical structure of constraint-shaped attractor spaces, not an
identity of substance.

\textbf{Measures.} The \textbf{Civilizational Stress Index (CSI)} is a
scalar reduction of the stress readout, measuring distance from the
stable regime; its construction and the correct interpretation of that
distance are specified in Document 2. The \textbf{Epistemic Coherence
Index (ECI)} operationalizes the information axis and is the
discipline's priority metrological build (§5).

\textbf{Dynamics and the central diagnostic.} The configuration evolves
under a deterministic field plus stochastic shock; within a poliform it
circulates through the secular cycles the cliodynamics literature
documents, and it transitions when it crosses a basin boundary. The
signature output is \textbf{distance-to-bifurcation} --- how close a
society is to qualitative reorganization --- expressed always as a
probability over a horizon, never as a date.

\textbf{Three ideas held as research hypotheses, not results.} The
discipline is candid that its most ambitious theoretical moves are not
yet established. \emph{Reflexivity} --- that a public model of social
dynamics becomes part of the system it models --- is treated as a
coupling to be characterized, not a solved problem. Candidate
\emph{quasi-conservation laws} (for legitimacy, coercive capacity,
information entropy) are \textbf{empirical hypotheses with stated
falsification conditions, not consequences of Noether's theorem}: the
discipline has no established action principle from which conservation
laws could be derived, and prior framing that implied otherwise is
corrected in Document 2 §8. The hypothesis of \emph{universality
classes} near transitions is exactly that --- a hypothesis, testable by
looking for scaling and data-collapse, not a property the framework may
assume.

\subsection{5. The measurement
commitment}\label{the-measurement-commitment}

A theory of civilizational dynamics is only as serious as its
measurement program, and the discipline binds itself to operational,
reproducible, openly published measurement. Every index is to be
published with its exact construction, its baseline and the defense of
that baseline, the data vintage, the normalization, the uncertainty, and
the code (Document 4 specifies this in full).

Two commitments deserve emphasis here because they are matters of
integrity, not only of method. First, \textbf{estimated quantities are
never to be presented as computed.} Much of the corpus today contains
illustrative estimates; these must be labeled as such in every artifact,
public or internal, with no computed-looking precision standing in for
measurement that has not been done. Second, \textbf{the ECI is the
highest-priority empirical investment of the program.} The information
layer is the discipline's distinctive contribution, and it is also its
least mature measure; credibility depends on \emph{building} validated
epistemic-coherence measurement --- a standing factual-consensus
instrument, cross-community media-overlap and network pipelines, and the
best feasible historical reconstruction --- rather than asserting ECI
values.

On the substance of the information layer, the discipline advances
specific \textbf{mechanisms}, stated as such and kept separate from
stronger claims of coordinated intent that would be harder to establish
and would pull toward unfalsifiability: that engagement-optimizing
platforms systematically amplify emotionally arousing,
identity-confirming, low-accuracy content, degrading the shared factual
commons; that above a threshold of fragmentation populations undergo an
\textbf{epistemic phase separation} into near-hermetic communities that
diverge on basic factual premises, disabling the self-correction that
democratic competition presupposes; that a \textbf{velocity mismatch}
between institutional and informational timescales prevents correction
from operating faster than destabilization; and that the dominant
information infrastructure is \textbf{structurally misaligned} with
social stability because its revenue depends on maximizing arousal.
These are the engine of the epistemic-layer research program (Documents
4 and 5).

\subsection{6. Validation architecture}\label{validation-architecture}

This section is the heart of the discipline's claim to be a science. The
standards below are binding and stated in advance.

\textbf{Retrodiction.} A minimum historical test set must be correctly
characterized --- meaning the framework must show elevated stress
preceding these transitions --- or the framework is falsified and must
be revised or abandoned: the late Roman Republic, the Weimar Republic,
the French Revolution, the Soviet collapse, the 1848 European
revolutions, and the United States 2008--2024. The success bar is
correct characterization in at least 80\% of the set. The worked
retrodictions live in \textbf{Document 6}, each tagged to the
evidentiary tier the available data support (the distinction between
locating attractors and reconstructing dynamics is drawn in Document 2
§3).

\textbf{Prospective prediction.} All predictions are probability
distributions over defined outcomes with explicit update conditions and
resolution criteria, timestamped and published before the window opens,
and scored with Brier scores upon resolution. The discipline targets
calibration exceeding the best superforecasters on structural questions
with three-to-ten-year horizons --- a high bar, achievable only if the
theory is sound. The public, auditable record is the mechanism of
accountability (§9).

\textbf{Falsifiability conditions.} The core claims fail if: the stress
index, constructed prospectively with pre-specified weights,
systematically does not precede transitions in the test set; prospective
predictions, tracked over a ten-year window, fail to beat a naive
baseline on Brier score; cross-civilizational analysis reveals no common
structural variables predicting transitions across the set; or the
information/epistemic layer, controlled for economic and demographic
variables, shows no independent contribution to disruption probability.
Each is a way the discipline could be wrong, stated so that being wrong
is possible.

\subsection{7. Relationship to existing
disciplines}\label{relationship-to-existing-disciplines}

Civilodynamics is integrative by design and collaborative by
disposition. The table states what it borrows, where it departs, and
where the most productive collaboration lies. The earlier ``supersedes''
framing has been retired in favor of ``departs from,'' which is both
more accurate and more conducive to the alliances the discipline needs.

{\def\LTcaptype{none} % do not increment counter
\begin{longtable}[]{@{}
  >{\raggedright\arraybackslash}p{(\linewidth - 6\tabcolsep) * \real{0.2182}}
  >{\raggedright\arraybackslash}p{(\linewidth - 6\tabcolsep) * \real{0.1636}}
  >{\raggedright\arraybackslash}p{(\linewidth - 6\tabcolsep) * \real{0.2545}}
  >{\raggedright\arraybackslash}p{(\linewidth - 6\tabcolsep) * \real{0.3636}}@{}}
\toprule\noalign{}
\begin{minipage}[b]{\linewidth}\raggedright
Discipline
\end{minipage} & \begin{minipage}[b]{\linewidth}\raggedright
Borrows
\end{minipage} & \begin{minipage}[b]{\linewidth}\raggedright
Departs from
\end{minipage} & \begin{minipage}[b]{\linewidth}\raggedright
Collaboration zone
\end{minipage} \\
\midrule\noalign{}
\endhead
\bottomrule\noalign{}
\endlastfoot
\textbf{Cliodynamics} & Structural-demographic model, secular cycles,
Seshat, prediction methodology & Thin information layer; primarily
historical calibration & \textbf{Direct collaboration; the parent
program and foundation} \\
Complexity science & Phase transitions, emergence, power laws, networks
& Metaphorical rather than literal application to societies & Core
mathematical language; SFI as institutional partner \\
Physics & Dynamical systems, statistical mechanics, universality (as
hypothesis) & Not departed from --- borrowed at the mathematical level &
Non-equilibrium statistical mechanics; sociophysics \\
Superforecasting & Probabilistic calibration, Brier scoring,
epistemic-humility norms & Atheoretical; short horizons & Calibration
methodology; prediction-tournament infrastructure \\
Economics & Mathematical formalism, behavioral economics, game theory &
Equilibrium assumptions; exogenous politics & Complexity economics;
behavioral foundations \\
Political science & Institutional typology, conflict datasets, electoral
studies & Event-level focus; structural-prediction failure & Comparative
institutional analysis; V-Dem \\
Sociology & Social-capital measures, institutional theory, surveys &
Static structural focus; prediction aversion & Institutional
measurement; micro-level data \\
Anthropology & Deep comparative data, Seshat, collapse studies &
Anti-quantitative tradition; scale incoherence & Cross-cultural
comparison; collapse studies (Tainter) \\
Cognitive science & Belief formation, motivated reasoning, identity
dynamics & Rarely connects to macro dynamics & Micro-foundations of the
epistemic layer \\
\end{longtable}
}

\subsection{8. The twenty-year arc}\label{the-twenty-year-arc}

The detailed, dependency-aware, continuously updated program is
\textbf{Document 8}; this is only its shape. The trajectory runs from
foundational theory and a validated measurement framework, through the
construction of historical and real-time databases and the first
cross-civilizational validation studies, to a formal dynamical treatment
with a public prediction platform, then to cross-civilizational testing
and theoretical refinement, and finally to full institutionalization ---
a journal, a curriculum, an institute, and a self-reproducing community
of researchers. The arc is calibrated against what comparable
communities have achieved: the Santa Fe Institute turned complexity
science from provocation into discipline in roughly two decades, and
cliodynamics moved quantitative history from the fringe to a recognized
program in a similar span. The governing principle across the arc is
that \textbf{rigor leads reach} --- priority is established early and
cheaply, but public promotion of predictive claims waits until the
measurement infrastructure is real.

\subsection{9. Institutional
architecture}\label{institutional-architecture}

The institutional vision is long-horizon and aspirational; its
sequencing is in Document 8. It comprises four elements.

\textbf{The Institute}, modeled on the Santa Fe Institute: an
independent, cross-disciplinary research center with no departmental
structure to enforce disciplinary orthodoxy, whose permanent members
span complex-systems physics, dynamical-systems mathematics,
quantitative history, complexity and behavioral economics, and
comparative social science, with no single discipline holding a
majority. Funding is to be diversified across foundations, government
research agencies, and international bodies to protect independence.

\textbf{The Journal of Civilodynamics}, publishing only work that is
quantitative, falsifiable, and cross-scale, operating open-access with
mandatory pre-registration for empirical studies.

\textbf{The Curriculum}, purpose-built across mathematics (dynamical
systems, stochastic processes, networks, Bayesian inference), physics
(statistical mechanics, phase-transition theory), empirical social
science (quantitative history, comparative analysis, survey
methodology), data science (pipelines, NLP, machine learning), and
philosophy of science (falsificationism, the epistemology of complex
systems, the ethics of predictive social science).

\textbf{The Public Prediction Platform}, which publishes every
structural prediction with its claim, probability, the state-vector
components it rests on, its update conditions, and its resolution
criteria, and scores all of them with Brier scores on an openly
accessible record. This is simultaneously the discipline's
accountability mechanism and its strongest credibility instrument ---
precisely because it holds the field to a standard most fields avoid.

\subsection{10. Ethical commitments}\label{ethical-commitments}

The ability to characterize civilizational states probabilistically
creates real obligations. The discipline binds itself to four.

\textbf{Transparency.} Predictions are published openly so that their
influence on outcomes --- the reflexivity effect --- is distributed
across the population rather than concentrated in the hands of any actor
positioned to benefit from or accelerate a predicted disruption.

\textbf{Funding integrity.} Civilodynamicists do not accept funding or
employment from actors whose interests are served by the destabilization
of the systems being modeled; the Institute is to maintain an ethics
board with real authority over publication and funding.

\textbf{Dual-use care.} The discipline develops explicit protocols for
findings that carry immediate crisis implications, analogous to those
governing the publication of dual-use biological research.

\textbf{Epistemic honesty as ethics.} The commitments of §5 --- never
presenting estimated values as computed, demoting contested claims to
operationalized and falsifiable sub-hypotheses, and publishing failures
alongside successes --- are treated not merely as methodology but as the
ethical core of a predictive social science, the safeguard against a
discipline that flatters its founders' priors or its funders' interests.

\subsection{11. Founding statement}\label{founding-statement}

Human civilizations are the most complex dynamical systems we know of,
exhibiting the full repertoire of complex-system behavior ---
self-organization, emergence, phase transitions, criticality, cascading
failure --- and they are the systems whose behavior most urgently
demands understanding, because the consequences of their transitions are
borne by billions. The existing sciences have produced deep knowledge of
the parts. What has been missing is a discipline willing to commit to
the whole: to treat the civilization as the unit of analysis, to model
the coupling between its subsystems, to build the measurement
infrastructure its state requires, and to hold itself to the
falsifiability standards that separate science from sophisticated
description.

Civilodynamics makes that commitment, with full awareness of the
difficulty. The reflexivity problem is real, the measurement challenges
are substantial, the formalism is partly undeveloped, and the historical
data are imperfect. These are research problems, not reasons to abandon
the work --- and the alternative, continuing to analyze civilizational
dynamics without a discipline designed to do so, has costs now
measurable in historical events.

\begin{quote}
\emph{The goal is not to be Hari Seldon. The goal is to build the
instrument that made Seldon possible --- to develop the measurement
theory, the mathematical formalism, and the institutional infrastructure
that a predictive science of civilizational dynamics requires.}
\end{quote}

The founding of this discipline begins here. The first commitments are
stated, the first falsifiability conditions set, and the corrected
foundations laid in the documents this charter governs. The clock starts
now.

\begin{center}\rule{0.5\linewidth}{0.5pt}\end{center}

\emph{Civilodynamics Initiative · Foundational Charter · Version 2.0 ·
June 2026. This charter is the constitutional document of the corpus;
Document 2 governs all mathematics, Document 8 the living program.
Changes to the discipline's core commitments are made here and
propagated forward.}

\end{document}
